% ----- CHAUPAI 4 -----

\newpage

\subsection*{\deva{चतुर्थ चौपाई} (Fourth Chaupai)}
\addcontentsline{toc}{subsection}{\deva{चतुर्थ चौपाई} (Fourth Chaupai) — \deva{कंचन बरन...}}

\begin{minipage}[t]{0.55\textwidth}
\vspace{0pt}

{\large \linespread{1.5}\selectfont
\deva{कंचन बरन बिराज सुबेसा।}\\
\deva{कानन कुंडल कुंचित केसा॥}
\par}

\vspace{0.5em}

{\large \linespread{1.5}\selectfont
\telu{కంచన బరన బిరాజ సుబేసా।}\\
\telu{కానన కుండల కుంచిత కేసా॥}
\par}

\vspace{0.5em}

{\large \linespread{1.3}\selectfont
\textit{kañcana barana birāja subesā |}\\
\textit{kānana kuṇḍala kuñcita kesā ||}
\par}
\end{minipage}%
\hfill
\begin{minipage}[t]{0.42\textwidth}
\vspace{0pt}
\begin{tcolorbox}[anchorbox]
\textbf{The Reward of Education:} Why the golden glow and beautiful earrings (\textit{Kanan Kundal})? According to tradition, when Hanuman completed his immense, relentless education under the Sun God, Surya was so impressed by his student that he gifted him a fraction of his own blinding brilliance and divine earrings as a graduation gift. True radiance is not just physical, it is the result of dedication, hard work and earned through good education.
\vspace{0.5em}\newline His brilliant, golden radiance was not inherited, but earned through the disciplined fire of his actions.\newline\null\hfill\textit{Kishkindha Kanda, 66}
\end{tcolorbox}
\end{minipage}

\vspace{0.5em}
\subsubsection*{\textbf{Visual Rhythm Map:}}
\begin{table}[h!]
\centering
\renewcommand{\arraystretch}{1.2}
\setlength{\tabcolsep}{0pt}
\begin{tabularx}{\textwidth}{|*{16}{>{\centering\arraybackslash\hsize=1.0\hsize}X|}}
\hline
\multicolumn{4}{|>{\centering\arraybackslash\hsize=4.0\hsize}X|}{\textbf{\guru{कं}\laghu{च}\laghu{न}} \par (4)} &
\multicolumn{3}{>{\centering\arraybackslash\hsize=3.0\hsize}X|}{\textbf{\laghu{ब}\laghu{र}\laghu{न}} \par (3)} &
\multicolumn{4}{>{\centering\arraybackslash\hsize=4.0\hsize}X|}{\textbf{\laghu{बि}\guru{रा}\laghu{ज}} \par (4)} &
\multicolumn{5}{>{\centering\arraybackslash\hsize=5.0\hsize}X|}{\textbf{\laghu{सु}\guru{बे}\guru{सा}} \par (5)} \\
\hline
\end{tabularx}
\vspace{-1.5em}
\end{table}

\begin{table}[h!]
\centering
\renewcommand{\arraystretch}{1.2}
\setlength{\tabcolsep}{0pt}
\begin{tabularx}{\textwidth}{|*{16}{>{\centering\arraybackslash\hsize=1.0\hsize}X|}}
\hline
\multicolumn{4}{|>{\centering\arraybackslash\hsize=4.0\hsize}X|}{\textbf{\guru{का}\laghu{न}\laghu{न}} \par (4)} &
\multicolumn{4}{>{\centering\arraybackslash\hsize=4.0\hsize}X|}{\textbf{\guru{कुं}\laghu{ड}\laghu{ल}} \par (4)} &
\multicolumn{4}{>{\centering\arraybackslash\hsize=4.0\hsize}X|}{\textbf{\guru{कुं}\laghu{चि}\laghu{त}} \par (4)} &
\multicolumn{4}{>{\centering\arraybackslash\hsize=4.0\hsize}X|}{\textbf{\guru{के}\guru{सा}} \par (4)} \\
\hline
\end{tabularx}
\vspace{-1.5em}
\end{table}

\vspace{1em}
\textbf{ \deva{सरल-संस्कृतम्}:}\\
 \deva{भवतः कञ्चन-वर्णः (सुवर्ण-वर्णः) अस्ति, भवान् सुन्दर-वेषेण च विराजते।}\\
 \deva{भवतः कर्णयोः कुण्डले स्तः, केशाः च कुञ्चिताः सन्ति॥}\\

Your complexion is golden, and you are adorned with beautiful attire.\\
You wear earrings on your ears, and your hair is curly.


\vspace{1em}
\newpage
\textbf{\deva{प्रतिपदार्थः} (Meaning of each word):}
\begin{table}[h!]
\centering
\renewcommand{\arraystretch}{1.2}
\begin{tabularx}{\textwidth}{@{} l l X X @{}}
\toprule
\textbf{Awadhi} & \textbf{Sanskrit} & \textbf{English} & \textbf{Grammar \& Notes} \\
\midrule
\deva{कंचन} / \telu{కంచన}  &  \deva{कञ्चन / सुवर्णम्}  &  Gold  &   \\
\deva{बरन} / \telu{బరన}  &  \deva{वर्णः}  &  Complexion / Color  & \\ 
\deva{बिराज} / \telu{బిరాజ}  &  \deva{विराजते}  &  Shines / Is adorned  & \\ 
\deva{सुबेसा} \gotchaicon / \telu{సుబేసా}  &  \deva{सुवेषः}  &  Beautiful attire  & \\ 
\deva{कानन} / \telu{కానన}  &  \deva{कर्णयोः}  &  In the ears  & \\ 
\deva{कुंडल} / \telu{కుండల}  &  \deva{कुण्डले}  &  Earrings  &   \\
\deva{कुंचित} / \telu{కుంచిత}  &  \deva{कुञ्चिताः}  &  Curly  &   \\
\deva{केसा} / \telu{కేసా}  &  \deva{केशाः}  &  Hair  & \\ 
\bottomrule
\end{tabularx}
\end{table}

\begin{tcolorbox}[gotchabox]
\begin{itemize}
    \item \textbf{Phonetic Shifts:} \deva{सुवेष} (Suvesha) becomes \deva{सुबेसा} (Subesaa) — exhibiting both the \deva{व} $\rightarrow$ \deva{ब} and \deva{श/ष} $\rightarrow$ \deva{स} shifts simultaneously! Likewise, \deva{केश} (Kesha) becomes \deva{केस} (Kesa).
\end{itemize}
\end{tcolorbox}



\newpage
